\clearpage\fancyhead[L]{\small COMET.}\hypertarget{native-348-lower}{}\label{native-348-lower}
\begin{center}\textbf{COMET.}\par\vspace{4pt}\hrule\vspace{4pt}\fontsize{10}{12}\selectfont Monatliche Correspondenz. Band XXIV. S. 595\,..\,598. 1811 December.\par\vspace{4pt}\hrule\vspace{4pt}\end{center}

\par\indent In Göttingen ward der Comet sogleich am 9. December, wo mein Brief dort einging, am Abend von den Herren \textsc{Gauss} und \textsc{Harding} aufgefunden. Am Kreismicrometer machte \textsc{Gauss} folgende Beobachtungen:

\begin{center}
\fontsize{10}{12}\selectfont\setlength{\tabcolsep}{4pt}\renewcommand{\arraystretch}{1.15}\begin{tabular}{l|c|c|c}
 & Mittlere Zeit & AR. Comet & Decl. Austr. \\
 & in Göttingen & & \\ \hline
Decemb. 9 & $10^{\mathrm u}\;\;6^{\mathrm m}\;52^{\mathrm s}$ & $63^\circ\;49'\;41''4$ & $10^\circ\;21'\;55''5$ \\
11 & $10\;\;34\;\;\;\;1$ & $63\;\;33\;\;18.0$ & $8\;\;39\;\;46.4$ \\
12 & $\;\;8\;\;\;\;5\;\;52$ & $63\;\;26\;\;25.3$ & $7\;\;54\;\;25.9$ \\
\end{tabular}
\end{center}

\SourceLeaf{349}
\par\indent „Hier, schrieb uns \textsc{Gauss}, Göttingen, am 15. Dec. 1811, auch meine vorläufigen Elemente, die freilich nur als eine Annäherung anzusehen sind, aber doch zeigen, dass der Comet ein \emph{neuer} ist, und an Licht nicht mehr zunehmen wird.

\Needspace{16\baselineskip}
\begin{center}\fontsize{10}{12}\selectfont\setlength{\tabcolsep}{4pt}\renewcommand{\arraystretch}{1.15}\begin{tabular}{l}
Zeit des Periheliums 1811. Nov. 12,6225 Gött. Mer.\\
\end{tabular}\end{center}
\begin{center}\fontsize{10}{12}\selectfont\setlength{\tabcolsep}{3pt}\renewcommand{\arraystretch}{1.15}\begin{tabular}{llr}
Länge des Perihels & .\;.\;.\;.\;.\;.\;. & $48^\circ\;30'\;20''$\\
Log. des perih. Abstandes & .\;.\;.\;.\;. & $0.20160$\\
Aufsteigender Knoten & .\;.\;.\;.\;.\;. & $92^\circ\;46'\;59''$\\
Neigung der Bahn & .\;.\;.\;.\;.\;.\;. & $31\;37\;55$\\
\end{tabular}\end{center}
\noindent Bewegung rechtläufig


\par\indent „Der Comet kann noch \emph{sehr} lange sichtbar sein, wenn seine Lichtstärke es nicht hindert. Ich finde 1812 Januar 31 $10^{\mathrm u}$ AR $67\frac{1}{2}^\circ$ Decl. $25\frac{1}{2}^\circ$ Nördl. Lichtstärke 0.178.

\begin{center}
\begin{tabular}{l}
Lichtstärke 1811 November 18 $\;=\;0.713$ \\
\qquad --- \qquad --- December 12 $\;=\;0.698$ \\
\end{tabular}
\end{center}


\Needspace{20\baselineskip}\par\indent Herr \textsc{Gerling}, der mit diesen Elementen alle seither bekannt gewordenen Beobachtungen verglich, fand folgende Resultate:

\begin{center}
\fontsize{10}{12}\selectfont\setlength{\tabcolsep}{6pt}\renewcommand{\arraystretch}{1.15}\begin{tabular}{l|c|c|l}
 & \multicolumn{2}{c|}{Abweichung} & \\
 & in AR. & in Decl. & Beoobachter \\ 
\hline
Novemb. 18 & $+32''$ & $-101''$ & v. Zach \\
19 & $+80$ & $+46$ & v. Zach \\
20 & $+23$ & $+4$ & v. Zach \\
21 & $+3$ & $+15$ & v. Zach \\
Decemb. 8 & $+19$ & .\;.\;.\;. & v. Lindenau \\
9 & $+11$ & $+12$ & v. Lindenau \\
9 & $+8$ & $-55$ & Olbers \\
9 & $+1$ & $+1$ & Gauss \\
11 & $-10$ & $-16$ & Gauss \\
12 & $+1$ & $+7$ & Gauss \\
\end{tabular}
\end{center}


\par\indent Um die künftigen Beobachtungen des Cometen zu erleichtern, berechnete Herr \textsc{Nicolai} nach obigen Elementen folgende Ephemeride:

\begin{center}
\fontsize{10}{12}\selectfont\setlength{\tabcolsep}{4pt}\renewcommand{\arraystretch}{1.15}\begin{tabular}{l|c|c|c|c}
Mitternacht & AR. & Declinat. & \shortstack{Log. dist.\\$a\;\EarthSymbol$} & \shortstack{Licht-\\stärke} \\ \hline
1811 December 21 & $62^\circ\;36'$ & $0^\circ\;7$ S. & 9.8879 & 0.596 \\
29 & $62\;\;25$ & $6\;\;20$ N. & 9.9185 & 0.496 \\
1812 Januar 6 & $62\;\;49$ & $12\;\;11$ --- & 9.9562 & 0.397 \\
14 & $63\;\;47$ & $17\;\;15$ --- & 9.9984 & 0.310 \\
22 & $65\;\;17$ & $21\;\;33$ --- & 0.0427 & 0.239 \\
30 & $67\;\;17$ & $25\;\;\;\,8$ --- & 0.0872 & 0.183 \\
Februar 7 & $69\;\;41$ & $28\;\;\;\,7$ --- & 0.1310 & 0.141 \\
\end{tabular}
\end{center}


\par\indent Hiernach nimmt die Lichtstärke des Cometen schon wieder sehr merklich ab, und da der Mondschein nun eine Zeitlang seine Sichtbarkeit hindern wird, so bleibt es wohl noch zweifelhaft, ob es gelingen wird, noch eine Reihe guter Beobachtungen dieses Cometen zu erhalten. Allein die schöne Übereinstimmung mit allen Beobachtungen zeigt, dass schon jetzt sein Lauf durch obige Elemente sehr genähert bestimmt ist, so dass er bei einer dereinstigen Wiedererscheinung von unsern Nachkommen ohne Mühe wieder erkannt werden kann.
